% =============================================================================
% Data and construction of the measures
% Milano housing project -- FSS group paper
% Source: docs/paper/sezione_data_e_misure.md (same content, LaTeX version)
% Standalone: compile with pdflatex (twice). To include in the main paper,
% copy everything between \begin{document} and \end{document}.
% =============================================================================
\documentclass[11pt,a4paper]{article}

\usepackage[utf8]{inputenc}
\usepackage[T1]{fontenc}
\usepackage{lmodern}
\usepackage[margin=2.5cm]{geometry}
\usepackage{amsmath,amssymb}
\usepackage{booktabs}
\usepackage{tabularx}
\usepackage{array}
\usepackage{pdflscape}
\usepackage[authoryear,round]{natbib}
\usepackage[hidelinks]{hyperref}
\usepackage{xurl}

\newcolumntype{L}{>{\raggedright\arraybackslash}X}
\newcommand{\eur}{EUR}

\begin{document}

% -----------------------------------------------------------------------------
\section{Data}\label{sec:data}

\paragraph{Rents.} Housing costs are measured with the quotations of the
\emph{Osservatorio del Mercato Immobiliare} (OMI) of the Italian Revenue
Agency, distributed by the Comune di Milano open-data portal (dataset ds1996,
first semester 2004 to second semester 2025). For each homogeneous zone,
property type and state of repair, OMI publishes a minimum--maximum interval
of market and rental values per m$^2$ \citep{omi}; in ds1996 rents are in
\eur{} per m$^2$ of gross floor area per month and sale values in \eur{} per
m$^2$. These intervals describe normal market values in the zone; they are
neither averages of registered contracts nor transaction prices. We keep
``civil dwellings'' (type 20) in ``normal'' condition, one quotation per zone
and semester, and use the midpoint of the range. The number of quoted zones
falls from 55 (2004--2013) to 38 (from 2014; 39 from 2018) and rises to 40
after the 2024 revision (Appendix~\ref{app:omi}).

\paragraph{Income.} Dataset ds531 reports, for the whole municipality, the
number of taxpayers (\emph{persone fisiche}) and the amount of \emph{reddito
complessivo} in eight income classes, for tax years 2008--2024. Tax years
2004--2007 come from the municipal files of the Ministry of Economy and
Finance (MEF), which use the same classes. \emph{Reddito complessivo} is an
individual taxpayer's total income before deductions and taxes. It is not
household income, and it is not available below the city level. Tax year $t$
(income earned in $t$, declared in $t+1$) is matched to rents of calendar year
$t$. Postcode-level MEF data are used only in CAP-level analyses.

\paragraph{Population and moves.} DS205 provides residents, births, deaths,
emigrants and immigrants for each \emph{Nucleo di Identit\`a Locale} (NIL) and
year, 2011--2024. Four further datasets split registry cancellations and
registrations into three reasons for 2004--2024: another Italian
municipality, abroad, and ``other'' (administrative removals such as
untraceable residents). Dataset ds27 gives residents by NIL and single year of
age for 1999--2025. DS205 and ds27 are stocks; the movement data are flows,
which never record moves between NILs.

\paragraph{Metro stations.} Dataset ds535 (municipal public-transport stops)
gives the name, lines and coordinates of each metro station, 21 of them on
line M4. Opening dates of the three M4 sections are added by hand
(26 November 2022, 4 July 2023 and 12 October 2024).

\paragraph{Analytical sample.} The unit of observation is the NIL-year, on the
88 current NILs. We drop eight NILs with fewer than 500 residents in 2019
(parks, rail yards). Of the remaining 80, the main panel keeps the 74 NILs
whose rent is observed in every year from 2004 to 2024. This balanced panel
contains 1,332 NIL-years for the main index (2007--2024) and 1,554 for the
gross-income control (2004--2024). Income is a single city-wide series. For
2004--2010, residents are set equal to their 2011 values.

% -----------------------------------------------------------------------------
\section{Construction of the measures}\label{sec:measures}

\subsection{From OMI zones to NILs}\label{sec:xwalk}

OMI zones and NILs do not nest. We therefore overlay a regular grid of points
(about $27\times 28$~m) on the city and assign each point to a NIL and to the
OMI zone in force under each zoning scheme $k$. Point counts approximate
intersection areas. For NIL $i$ and semester $s$ with scheme $k(s)$, the rent
is
\begin{equation}\label{eq:rent}
R_{is}=\frac{\sum_{z\in Q_s} n^{k(s)}_{iz}\,R_{zs}}{\sum_{z\in Q_s} n^{k(s)}_{iz}},
\qquad
c_{is}=\frac{\sum_{z\in Q_s} n^{k(s)}_{iz}}{N_i},
\end{equation}
where $R_{zs}$ is the OMI midpoint, $n^{k}_{iz}$ is the number of points of
NIL $i$ in zone $z$, $Q_s$ is the set of zones quoted in $s$, and $N_i$ is the
total number of points in the NIL. $R_{is}$ is set to missing when coverage
$c_{is}<0.80$. The annual rent $R_{it}$ is the mean of the available
semesters. Rents are intensive (\eur/m$^2$): the weights average them and
never allocate totals. The underlying assumption is that the quotation is
uniform within a zone and that every part of the NIL counts in proportion to
its area, so parks and rail yards weigh as much as residential blocks. NIL
rents are thus estimates.

\subsection{OMI zoning changes}\label{sec:zoning}

Each semester is linked to the zoning in force at the time: the 2011 polygons
up to 2013; for 2014S2--2024S1, the 2025 polygons with the zones split in 2024
recoded to their parents (D37+D38$\to$D15, D39+D40$\to$D30); and the 2025
polygons as published from 2024S2. The first semester of 2014 is a transition
period, with half of the zones still carrying old codes, and is dropped.
Across the 2013--2014 break, measured changes may partly reflect redrawn
boundaries.

\subsection{Median income from grouped data}\label{sec:median}

Let classes $j=1,\dots,J$ have bounds $[L_j,U_j)$ and frequencies $f_j$. Let
$N_t=\sum_j f_j$ count all taxpayers, including those with zero or negative
income, and let $F_j$ be the cumulative frequency. The median class $m$
satisfies $F_{m-1}<N_t/2\le F_m$. Assuming taxpayers are uniformly spread
within it,
\begin{equation}\label{eq:median}
Y^{med}_t = L_m+\frac{N_t/2-F_{m-1}}{f_m}\,(U_m-L_m).
\end{equation}
The median always lies in the \eur{} 15,000--26,000 class, so the open top
class never matters. The distribution is city-wide, so no territorial medians
are averaged. $Y^{med}_t$ is an interpolated estimate for the median
taxpayer, not a microdata median. Net income replaces the class bounds with
their after-tax values, $g_t(L_m)$ and $g_t(U_m)$. The function $g_t$ applies
the national income tax (IRPEF), employee tax credits, the regional and
municipal surcharges and the payroll bonuses in force in year $t$:
\begin{equation}\label{eq:netmedian}
Y^{net}_t = g_t(L_m)+\frac{N_t/2-F_{m-1}}{f_m}\,\bigl[g_t(U_m)-g_t(L_m)\bigr].
\end{equation}

\subsection{Affordability index}\label{sec:index}

The main index is the annual rent of a standard 60~m$^2$ dwelling in NIL $i$
divided by the city's median net income:
\begin{equation}\label{eq:index}
A_{it}=\frac{12\cdot 60\cdot R_{it}}{Y^{net}_t},\qquad
D_{it}=\mathbf{1}\left[A_{it}>0.40\right].
\end{equation}
$A_{it}$ is a ratio, reported in percent; a higher value means lower
affordability. The denominator is the same for every NIL, so $A_{it}$
measures how much of the typical Milanese taxpayer's net income a standard
rental in each NIL would absorb, and its variation across NILs comes only
from rents. The 40\% threshold is the OECD housing cost overburden standard:
the share of households spending ``more than 40\% of their disposable income
on housing costs (in line with Eurostat methodology)'', where indicator
HC1.2.3 uses the narrow definition of rent and mortgage costs
\citep{oecd2024,eurostat}. Our numerator, rent only, matches that narrow
definition and our denominator is after tax, but two adaptations remain. The
OECD ratio uses each household's actual costs and equivalised household
disposable income; $A_{it}$ uses a standardised market rent and one
individual's estimated net income. A gross-income version,
$A^{g}_{it}=12\cdot60\cdot R_{it}/Y^{med}_t$ with the 30\% rule of thumb, is
kept only as a robustness check from 2004. The headline statistic is the
share of residents living in NILs above the threshold,
\begin{equation}\label{eq:share}
S_t=\frac{\sum_i P_{it}D_{it}}{\sum_i P_{it}},
\end{equation}
computed on the balanced panel. It says where a median earner could rent a
standard flat at normal market rates, not how many households are actually
overburdened. Sale values enter no price-to-income ratio and no threshold.

\subsection{Exits and net migration by age}\label{sec:mobility}

Out-migration is $r^{ac}_{it}=1000\cdot C^{ac}_{it}/P_{i,t-1}$, where
$C^{ac}$ counts cancellations to another Italian municipality. This excludes
moves abroad and administrative removals, which account for 42--46\% of
cancellations in 2013--2014. Registry flows ignore entries, as well as moves
within Milan. For age group $c$ we therefore follow cohorts in the ds27
stocks:
\begin{equation}\label{eq:netmig}
NM^{c}_{it}=1000\cdot\frac{\sum_{a\in c}P_{i,a+1,t}-\sum_{a\in c}P_{i,a,t-1}}
{\sum_{a\in c}P_{i,a,t-1}},
\qquad c\in\{0\text{--}9,\;20\text{--}29,\;30\text{--}39,\;40\text{--}64\}.
\end{equation}
$NM^{c}_{it}$ is the net change in cohort size. It nets out inflows against
outflows, including moves between NILs. Births do not enter it, but it also
absorbs deaths, which are small below age 65, and registry corrections. It is
therefore net migration only in an approximate sense. A balance near zero may
hide large gross flows. Neither measure separates voluntary mobility from
involuntary displacement.

% -----------------------------------------------------------------------------
\begin{thebibliography}{9}

\bibitem[{Agenzia delle Entrate}(n.d.)]{omi}
Agenzia delle Entrate (n.d.).
\newblock \emph{Quotazioni immobiliari}. Osservatorio del Mercato Immobiliare.
\newblock \url{https://www.agenziaentrate.gov.it/portale/web/guest/schede/fabbricatiterreni/omi/banche-dati/quotazioni-immobiliari}
(accessed 25 November 2026).

\bibitem[{Eurostat}(n.d.)]{eurostat}
Eurostat (n.d.).
\newblock \emph{Glossary: Housing cost overburden rate}. Statistics Explained.
\newblock \url{https://ec.europa.eu/eurostat/statistics-explained/index.php?title=Glossary:Housing_cost_overburden_rate}
(accessed 25 November 2026).

\bibitem[{OECD}(2024)]{oecd2024}
OECD (2024).
\newblock \emph{HC1.2 Housing costs over income}. OECD Affordable Housing
Database, last updated 15 April 2024.
\newblock \url{https://www.oecd.org/content/dam/oecd/en/data/datasets/affordable-housing-database/hc1-2-housing-costs-over-income.pdf}
(accessed 25 November 2026).

\end{thebibliography}

\noindent\emph{Data:} Comune di Milano, Open Data portal, datasets ds1996,
ds531, DS205, ds1596, ds1600, ds1954, ds1956, ds27, ds535 and DS964; MEF,
IRPEF statistics by municipality (2004--2007) and by postcode (2019--2024).

% -----------------------------------------------------------------------------
\begin{landscape}
\begin{table}[p]
\centering
\caption{Data sources and coverage}\label{tab:sources}
\small
\begin{tabularx}{\linewidth}{@{}p{4.2cm}LLp{3.4cm}L@{}}
\toprule
Source (producer, dataset) & Content used & Original geography $\to$ final & Frequency, years & Use \\
\midrule
OMI quotations, Agenzia delle Entrate, via Comune di Milano ds1996 &
Min/max rent (\eur/m$^2$/month) and sale value (\eur/m$^2$), gross floor area; civil dwellings, normal condition; midpoint &
OMI zone (55 $\to$ 38/39 $\to$ 40) $\to$ NIL (area weights) &
Semester, 2004S1--2025S2 (2014S1 excluded) &
$R_{it}$, numerator of $A_{it}$; sale values for the M4 analysis \\
OMI zone polygons (2011 S1; 2025 S2) & Zone boundaries & -- & Two vintages & Spatial weights $n^k_{iz}$ \\
Comune di Milano ds531 & Taxpayers and amounts by 8 classes of \emph{reddito complessivo} & Municipality & Annual, tax years 2008--2024 & $Y^{med}_t$, $Y^{net}_t$ \\
MEF, IRPEF data by municipality & Same classes & Municipality & Annual, tax years 2004--2007 & $Y^{med}_t$ (net from 2007) \\
MEF, IRPEF data by postcode & Taxpayers, amounts, classes & 38 CAP & Annual, 2019--2024 & CAP-level analyses only \\
Comune di Milano DS205 & Residents, births, deaths, emigrants, immigrants, age shares & NIL (88) & Annual, 2011--2024 & $P_{it}$, weights, $r^{emig}_{it}$ \\
ds1600/ds1956, ds1596/ds1954 & Cancellations and registrations by reason, sex, citizenship & NIL (matched by name) & Annual, 2004--2024 & $r^{ac}_{it}$ \\
ds27 & Residents by single year of age and citizenship & NIL & Annual, 1999--2025 & $NM^c_{it}$, composition \\
NIL polygons (DS964) & Current NIL boundaries & -- & Current & Grid assignment \\
ATM stops, ds535 & Metro station coordinates and lines; M4 opening dates added by hand & Point $\to$ grid $\to$ NIL / OMI zone & Static (cohort by section) & M4 exposure \\
\bottomrule
\end{tabularx}
\end{table}
\end{landscape}

\begin{table}[htbp]
\centering
\caption{Descriptive statistics}\label{tab:descr}
\small
\begin{tabular}{@{}lllrrrrrr@{}}
\toprule
Variable & Unit & Years & $N$ & Mean & SD & Median & Min & Max \\
\midrule
OMI rent $R_{it}$ & \eur/m$^2$/month & 2007--24 & 1,332 & 9.08 & 3.07 & 8.20 & 5.26 & 25.11 \\
Affordability index $A_{it}$ & \% & 2007--24 & 1,332 & 36.99 & 11.95 & 33.96 & 22.14 & 94.79 \\
Gross control $A^g_{it}$ & \% & 2004--24 & 1,554 & 30.61 & 9.85 & 28.22 & 18.06 & 77.24 \\
$D_{it}$: $A_{it}>40\%$ & \% NIL-years & 2007--24 & 1,332 & 26.95 & -- & -- & -- & -- \\
$A^g_{it}>30\%$ & \% NIL-years & 2004--24 & 1,554 & 39.77 & -- & -- & -- & -- \\
Median gross income $Y^{med}_t$ & \eur & 2004--24 & 21 & 21,053 & 2,056 & 20,947 & 17,344 & 25,919 \\
Median net income $Y^{net}_t$ & \eur & 2007--24 & 18 & 17,670 & 1,663 & 17,415 & 15,590 & 21,559 \\
Taxpayers $N_t$ & persons & 2004--24 & 21 & 980,597 & 18,191 & 974,592 & 953,140 & 1,013,593 \\
Residents $P_{it}$ & persons & 2012--24 & 962 & 18,037 & 12,018 & 16,679 & 367 & 63,073 \\
Exits to other municipalities $r^{ac}_{it}$ & per 1,000 & 2012--24 & 962 & 23.97 & 12.36 & 21.52 & 8.57 & 172.68 \\
All emigrants $r^{emig}_{it}$ & per 1,000 & 2012--24 & 962 & 32.81 & 18.53 & 28.68 & 10.90 & 257.06 \\
Net migration, 0--9 & per 1,000 & 2012--24 & 962 & 4.49 & 48.06 & 0.00 & $-285.7$ & 750.0 \\
Net migration, 20--29 & per 1,000 & 2012--24 & 962 & 63.15 & 122.64 & 52.84 & $-376.5$ & 3,286.3 \\
Net migration, 30--39 & per 1,000 & 2012--24 & 962 & 19.48 & 52.46 & 14.55 & $-234.4$ & 754.1 \\
Net migration, 40--64 & per 1,000 & 2012--24 & 962 & $-3.75$ & 25.62 & $-5.11$ & $-166.7$ & 464.6 \\
\bottomrule
\end{tabular}
\begin{minipage}{\linewidth}
\footnotesize\vspace{4pt}
\emph{Notes:} balanced panel of 74 NILs; unweighted statistics (each
NIL-year counts once); no imputation. City-level income series have one
observation per year. Population-weighted share of residents in NILs with
$A_{it}>40\%$ ($S_t$): 45.4\% (2007), 26.8\% (2018), 42.9\% (2024). Extreme
values of $NM^c$ come from small cohorts in sparsely populated NILs.
\end{minipage}
\end{table}

% -----------------------------------------------------------------------------
\appendix

\section{OMI zoning changes and spatial harmonisation}\label{app:omi}

\paragraph{Grid.} Points are spaced $0.00035^\circ$ in longitude and
$0.00025^\circ$ in latitude (about $27\times28$~m) over the extent of the 38
CAP polygons, giving 239,111 points (about 181~km$^2$) inside the CAPs. Of
these, 0.28\% fall outside any NIL and 0.31\% outside the 2025 OMI zones. A
point on a shared boundary goes to the first polygon (164 points for the 2025
zones, 460 for the 2011 zones). Each point carries one NIL, one 2011 zone, one
2025 zone and its recoded 2014--2024 zone, so joining quotations to points
cannot duplicate observations. Coverage below 0.80 sets the rent to missing
rather than rescaling it. This is why six NILs (9, 34, 35, 39, 40, 59) leave
the balanced panel after 2014: zone C13 is quoted only in ``excellent''
condition, and some peripheral zones have no civil-dwelling quotation.

\begin{landscape}
\begin{table}[p]
\centering
\caption{OMI zoning changes and their treatment}\label{tab:zoning}
\small
\begin{tabularx}{\linewidth}{@{}p{2.3cm}LLp{3.6cm}L@{}}
\toprule
Period & Documented change (from the data) & Treatment in the project & Source & Comparability \\
\midrule
2004S1--2013S2 & Same 55 zone codes every semester & 2011 S1 polygons applied to all years & ds1996 codes; OMI 2011 S1 polygons & Assumes boundaries constant 2004--2013 (not verified) \\
2014S1 & Transition: 38 quoted zones, about half with old codes (B02--B06, C01--C09, D01--D09) and no matching polygon & Semester dropped; 2014 = 2014S2 only & ds1996 & 2014 annual value based on one semester \\
2014S2 & New zoning: zones renumbered and redrawn (e.g.\ B15--B21, C15--C20, D30--D36, E5--E8); 40 zones, 38 quoted for civil/normal & 2025 S2 polygons with D37+D38$\to$D15, D39+D40$\to$D30 & ds1996; OMI 2025 S2 polygons & Break between 2013 and 2014; no zone-to-zone correspondence imposed \\
2018S2 & C14 first quoted for civil dwellings in normal condition (39 zones) & Enters the weights when quoted & ds1996 & Coverage of affected NILs rises \\
2024S2 & D15 $\to$ D37 + D38; D30 $\to$ D39 + D40; Ortomercato moved from D15 to D16; D39 quoted only in ``excellent'' condition & 2025 S2 polygons as published; 2024 = mean of S1 (old scheme) and S2 (new) & ds1996; OMI 2025 S2 polygons & Recoded D15 polygon for 2014--2024 omits the Ortomercato area \\
\bottomrule
\end{tabularx}
\end{table}
\end{landscape}

\section{Variable dictionary}\label{app:dict}

\begin{table}[htbp]
\centering
\caption{Variable dictionary}\label{tab:dict}
\footnotesize
\begin{tabularx}{\linewidth}{@{}p{3.3cm}p{1.2cm}p{3.6cm}Lp{2cm}@{}}
\toprule
Paper name & Symbol & Code name (file) & Definition & Unit \\
\midrule
OMI rent (zone) & $R_{zs}$ & \texttt{R\_mid} (\texttt{omi\_civ\_sem}) & Midpoint of min/max rent, civil dwellings, normal & \eur/m$^2$/month \\
OMI sale value (zone) & $P_{zs}$ & \texttt{P\_mid} & Midpoint of min/max sale value & \eur/m$^2$ \\
NIL rent & $R_{it}$ & \texttt{R} (\texttt{aff\_nil\_year}) & Area-weighted, mean of semesters & \eur/m$^2$/month \\
OMI coverage & $c_{it}$ & \texttt{cov} & Minimum semester share of NIL area in quoted zones & share \\
Median gross income & $Y^{med}_t$ & \texttt{y\_med} (\texttt{milano\_income}) & Interpolated median \emph{reddito complessivo}, city & \eur/year \\
Median net income & $Y^{net}_t$ & \texttt{y\_med\_net} & Same, after IRPEF, surcharges and bonuses & \eur/year \\
Taxpayers & $N_t$ & \texttt{n\_contr} & All ds531 classes & persons \\
Affordability index & $A_{it}$ & \texttt{A60} & $12\cdot60\cdot R_{it}/Y^{net}_t$ & ratio \\
Above threshold & $D_{it}$ & \texttt{off\_A60} & $A_{it}>0.40$ & 0/1 \\
Gross control & $A^g_{it}$ & \texttt{A60g}, \texttt{off\_A60g} & $12\cdot60\cdot R_{it}/Y^{med}_t$; threshold 0.30 & ratio; 0/1 \\
Variants & -- & \texttt{A40}, \texttt{A60d} & 40~m$^2$ one income; 60~m$^2$ two incomes & ratio \\
Balanced panel & -- & \texttt{balanced} & Rent observed every year 2004--2024 & 0/1 \\
Share of residents above threshold & $S_t$ & \texttt{sh\_pop\_A60} & $\sum_i P_{it}D_{it}/\sum_i P_{it}$ & share \\
Residents & $P_{it}$ & \texttt{pop} (\texttt{ds205\_nil\_year}) & Registered residents (2004--10 = 2011) & persons \\
Exits to other municipalities & $r^{ac}_{it}$ & \texttt{r\_ac} & $1000\cdot C^{ac}_{it}/P_{i,t-1}$ & per 1,000 \\
Exits incl.\ abroad & -- & \texttt{r\_mov} & Same, including moves abroad & per 1,000 \\
All emigrants & $r^{emig}_{it}$ & \texttt{r\_emig} & Includes administrative removals & per 1,000 \\
Net migration by age & $NM^c_{it}$ & \texttt{nm\_c0009}, \texttt{nm\_c2029}, \texttt{nm\_c3039}, \texttt{nm\_c4064} & Cohort change per 1,000 of cohort in $t-1$ & per 1,000 \\
\bottomrule
\end{tabularx}
\end{table}

\end{document}
